\section{Conclusion}
\label{sec:conclusion}

We have presented GAWorld v2, a research stack for building, auditing, and
revising LLM-based artificial societies.  v2 extends the v1 platform along
five dimensions: a fourteen-plugin microkernel with a closed K1--K8
migration; multi-tenant, per-world, role-aware deployment; hub-and-node
distributed execution and playable cities; a Research Workbench that
compiles a paper or research idea into a pre-registered protocol and a
deterministic counterfactual verdict; and an Agent Playground of fourteen
game-based probes and an Agent Arena.  v2 also closes a long-standing gap
between documented and implemented behavior in the price-level module,
adds an indoor spatial tree, and persists organizations and governance as
a separate ledger from the per-step metrics.

The contribution of v2 is not a new cognitive architecture, a new
counterfactual estimator, or a new social theory.  It is a configurable
instrument whose plugins, hooks, and surfaces make the construction and
audit of an artificial-society experiment a tractable engineering task.
Every reported number in this paper is bound to a hash, an evidence
class, and the producing code path; every failed arm is reported as a
failure, not as a partial finding.  The five claim levels that v1
introduced are extended with a sixth (participant evaluation) and remain
the boundary against which substantive social claims are held.  GAWorld
v2 is offered as a substrate for the research community, not as a
substitute for it: the construction, the audit, and the revision of
artificial-society hypotheses remain the work of the researcher, and the
platform is only as good as the questions it is asked to make auditable.

Three directions are particularly open.  First, a more comprehensive
\emph{individual-consistency} evaluation that uses repeated standardized
probes before and after controlled memory perturbations, with a
pre-registered analysis.  Second, a more comprehensive
\emph{macro-correspondence} evaluation that maps a small number of
simulated constructs to external measures, with a separate calibration
sample and a separate testing sample.  Third, a more comprehensive
\emph{participant-evaluation} evaluation that uses the Track B scorecard
to test the three pre-registered laws under sufficient sample sizes,
and adds two more laws (rumor verification; referendum partisanship)
under the same gate.  These are the research steps that v2 makes
tractable and that v3, if it is written, will report on.

The platform's source code, the artifact ledger, the per-arm
completion audits, the Track A and Track B scorecards, the Research
Workbench pre-registered protocols, and the parallel-worlds counterfactual
results are versioned together; a reader who wants to inspect a claim can
follow the hash, the evidence class, and the world and node identifiers
back to the producing code path.  This is the level of audit that we
believe the field of LLM-based social simulation should aspire to, and
the level at which GAWorld v2 is offered.
